\chapter{Conclusões e Trabalhos Futuros}
\label{chap:conclusoes-e-trabalhos-futuros}

\section{Conclusões}
\label{sec:conclusoes}

Este trabalho teve como objetivo desenvolver uma ferramenta web capaz de interpretar modelos estruturais elaborados no FTOOL e auxiliar na análise e no pré-dimensionamento de pontes treliçadas de palitos de picolé, com aplicação voltada à Competição de Pontes com Palitos de Picolé (CP3) da Universidade Federal Rural do Semi-Árido. A solução desenvolvida integrou etapas anteriormente executadas de forma separada, desde a interpretação do arquivo estrutural até a apresentação das estimativas de material necessárias à composição das barras.

Em relação à interpretação dos modelos, foi implementado um módulo capaz de realizar a leitura de arquivos FTL, identificar as informações estruturais e reconstruir a geometria e a conectividade dos elementos. Os testes realizados demonstraram a consistência desse procedimento para os casos avaliados, incluindo o modelo utilizado na apresentação dos resultados, composto por 10 nós e 19 barras. Também foi verificada a equivalência entre diferentes formas de leitura do arquivo, permitindo seu processamento diretamente em memória na aplicação web.

A integração com a biblioteca anaStruct possibilitou a execução da análise estrutural a partir das informações interpretadas do FTOOL, com a obtenção de esforços internos, deslocamentos e reações de apoio. Na comparação dos esforços axiais do modelo analisado, submetido a um carregamento vertical de 2,00 kN em uma de suas faces, os resultados apresentaram diferenças absolutas não superiores a 0,001 N em relação aos valores exibidos pelo FTOOL. Essa correspondência constitui uma evidência da consistência da conversão e da análise numérica para o caso estudado.

Quanto ao pré-dimensionamento, foram implementadas rotinas para avaliar barras submetidas à tração e à compressão, considerando critérios de resistência, estabilidade e composição construtiva. A ferramenta permite determinar a quantidade estimada de camadas de cada elemento e contabilizar os palitos comerciais necessários ao longo de seu comprimento.

No modelo utilizado para a avaliação, o pré-dimensionamento resultou em uma soma de 154 camadas distribuídas entre as 19 barras, com estimativa de 430 palitos para uma face da ponte. Considerando duas faces idênticas, a estimativa correspondeu a 860 palitos, sem incluir elementos adicionais de ligação transversal ou perdas de fabricação. Esses resultados demonstraram a capacidade da implementação de produzir estimativas individualizadas de material a partir dos esforços obtidos na análise estrutural.

A integração das rotinas em uma aplicação web permitiu reunir a importação do arquivo FTL, a visualização preliminar da geometria, a configuração dos parâmetros de cálculo e a consulta dos resultados em um único ambiente. A utilização do Pyodide possibilitou a execução das rotinas Python diretamente no navegador, enquanto o emprego de um Web Worker permitiu separar o processamento computacional das operações da interface.

A comparação com uma planilha eletrônica preexistente evidenciou diferenças entre as formulações de pré-dimensionamento e os procedimentos de contabilização do material. Mesmo após a adoção do mesmo comprimento comercial de palito, as estimativas apresentaram diferenças expressivas, associadas aos critérios e às hipóteses empregados em cada abordagem. Esses resultados ressaltam a importância de explicitar os parâmetros e as formulações utilizados nas estimativas, não sendo possível interpretar a diferença numérica, isoladamente, como evidência de maior precisão de um dos procedimentos.

Apesar dos resultados obtidos, a ferramenta apresenta limitações relacionadas à modelagem bidimensional, às propriedades mecânicas adotadas para os palitos e à representação simplificada das seções compostas. O pré-dimensionamento considera a atuação solidária entre as camadas e não contempla integralmente os efeitos das ligações, das emendas e das imperfeições construtivas. Além disso, não foi realizada uma campanha experimental destinada à validação das estimativas de resistência ou à previsão da carga de ruptura das pontes construídas.

Dessa forma, considera-se que o objetivo geral do trabalho foi alcançado no escopo proposto, uma vez que foi desenvolvida uma ferramenta web funcional capaz de interpretar modelos do FTOOL, executar a análise estrutural e realizar o pré-dimensionamento de barras de palitos. A solução contribui para a integração das etapas computacionais envolvidas no projeto dessas estruturas e disponibiliza um recurso de apoio à análise e à estimativa de material no contexto da CP3.

\section{Trabalhos futuros}
\label{sec:trabalhos-futuros}

Como continuidade deste trabalho, recomenda-se a realização de estudos experimentais destinados à caracterização das propriedades mecânicas dos palitos de picolé e das ligações utilizadas na construção das barras. Ensaios de tração, compressão e flambagem poderiam fornecer parâmetros mais representativos do material e permitir avaliar a adequação das hipóteses empregadas no pré-dimensionamento.

Outra possibilidade consiste no aperfeiçoamento da representação das seções compostas, considerando a influência das emendas, da rigidez do adesivo e de possíveis deslocamentos relativos entre as camadas. Essa evolução permitiria investigar o comportamento dos elementos construídos com palitos colados de maneira mais próxima às condições reais de fabricação.

Também se propõe a ampliação da análise para modelos tridimensionais, de modo a considerar simultaneamente as duas faces da ponte, os elementos de travamento transversal e a distribuição dos carregamentos entre os componentes da estrutura. Essa ampliação possibilitaria avaliar efeitos que não são representados integralmente pelos modelos planos utilizados neste trabalho.

No âmbito computacional, a ferramenta poderia ser aprimorada para contemplar outros elementos e funcionalidades do formato FTL, além de ampliar as rotinas de verificação e os recursos de visualização dos resultados. A inclusão de mecanismos para comparação de diferentes configurações estruturais também permitiria apoiar a avaliação de alternativas de projeto.

Por fim, recomenda-se a realização de estudos comparativos envolvendo pontes construídas e ensaiadas experimentalmente a partir de configurações previamente analisadas pela ferramenta. A comparação entre as estimativas computacionais, as quantidades efetivamente utilizadas e os modos de falha observados possibilitaria identificar limitações do modelo e orientar o aperfeiçoamento dos critérios de pré-dimensionamento.
